\chapter{Metodología de extracción y formalización de criterios curriculares}\label{anexo:extraccion}

Este anexo documenta cómo se estructuró la base de conocimiento de la oferta académica del IPN en la Ciudad de México a partir de fuentes oficiales.

\section{Corpus curricular de nivel superior del IPN}
Se reunió un documento oficial por cada carrera de nivel superior que imparte el IPN en la Ciudad de México (66 documentos). De cada uno se procesaron tres componentes:
\begin{enumerate}
    \item \textbf{Objetivo general del programa:} define la misión formativa, el campo de aplicación y la orientación disciplinar.
    \item \textbf{Perfil de ingreso:} especifica los conocimientos previos, aptitudes y disposiciones esperadas en los aspirantes.
    \item \textbf{Perfil de egreso:} detalla las competencias profesionales, destrezas técnicas y campos laborales del graduado.
\end{enumerate}
Los documentos PDF se convirtieron a Markdown antes de enviarlos al modelo de lenguaje, lo que reduce el número de tokens.

\section{Directrices psicológicas para la formalización RIASEC}
La Licenciada en Psicología Paola Alejandra Luna Cabrera estableció las reglas de correspondencia entre los descriptores curriculares y la taxonomía RIASEC~\cite{Holland1997}. Los criterios generados por el modelo de lenguaje se revisaron con esas reglas (esquema \textit{Human-in-the-Loop}):
\begin{itemize}
    \item \textbf{Realista ($R$):} manipulación de instrumentos y maquinaria, diseño mecánico, circuitos físicos, pruebas de laboratorio y actividades operativas.
    \item \textbf{Investigador ($I$):} habilidades analíticas, resolución matemática, modelado formal, diseño de algoritmos, análisis biomédico e investigación científica.
    \item \textbf{Artístico ($A$):} innovación visual, ergonomía, diseño de interfaces, síntesis creativa y comunicación conceptual. El verbo ``diseñar'' no implica por sí mismo esta dimensión.
    \item \textbf{Social ($S$):} trabajo multidisciplinario, comunicación interpersonal, docencia, bioética, administración de personal y atención comunitaria.
    \item \textbf{Emprendedor ($E$):} liderazgo de proyectos tecnológicos, gestión de costos, negociación, creación de empresas de base tecnológica y toma de decisiones estratégicas.
    \item \textbf{Convencional ($C$):} cumplimiento de normas y estándares, auditoría, control de calidad, seguridad informática y gestión de registros.
\end{itemize}

\section{Estructura del prompt de extracción}
La extracción usó Gemini 2.5 Pro con temperatura 0.2, una carrera por llamada, y respuesta estructurada. La figura~\ref{fig:prompt-extraccion} resume las instrucciones.

\begin{figure}[htbp]
\centering
\begin{minipage}{0.95\textwidth}
\ttfamily\small
\textbf{INSTRUCCIONES DEL PROMPT DE EXTRACCIÓN (RESUMEN)}\par\smallskip
\textbf{Rol:} experto en psicología educativa y orientación vocacional (modelo Holland).\par
\textbf{Entrada:} documento de la carrera (objetivo, perfil de ingreso y perfil de egreso).\par\smallskip
\textbf{Reglas:}\par
1. Un criterio por cada dimensión R, I, S, E y C. Cero o uno para A.\par
2. Cada criterio es una habilidad: verbo en infinitivo, objeto y contexto. Sin sujeto, sin adornos, sin duplicados.\par
3. Resumen de 2 a 8 palabras (máximo 12).\par
4. Códigos de dimensión: R=1, I=2, A=3, S=4, E=5, C=6.\par
5. Orden de fuentes: PERFIL, luego INTERNET y por último CONTINGENCIA. A nunca se obtiene por contingencia.\par
6. En INTERNET no se aceptan blogs, foros, wikis ni respuestas de IA generativa.\par\smallskip
\textbf{Salida:} tres hojas: criterios, fuentes y cobertura\_riasec.
\end{minipage}
\caption{Resumen de las instrucciones del prompt de extracción}
\label{fig:prompt-extraccion}
\end{figure}

\section{Persistencia en la base de datos}
Los criterios validados se insertaron en la tabla \texttt{criterios} de PostgreSQL, vinculados a las tablas \texttt{carreras}, \texttt{tipos\_criterio} y \texttt{tipos\_riasec}. La carga verifica antes de escribir que existan 66 carreras, 331 criterios de perfil y 66 principales, y que no haya claves duplicadas. El origen de cada criterio (perfil, internet o contingencia) se guarda en el campo \texttt{type\_extraction}. Estos criterios sirvieron para generar los aspirantes sintéticos y para entrenar los modelos Random Forest.

\endinput
